\begin{proposition}[B.4.10 (Proved): an explicit parabolic family and its nonextending exceptional correction]
Fix an integer $n\ge1$ and holomorphic units $a(s),\kappa(s)$ on a
disc, with $0<|a(s)|<1$. After shrinking the disc there is a proper
holomorphic submersion $\mathcal K\to\Delta_s$, with an embedded smooth
proper elliptic family
\[
 \mathcal E=(\C^*\times\Delta_s)/
       ((v,s)\sim(a(s)v,s)).
\]
Its central surface is a parabolic Inoue surface $S_n$, with
$\pi_1(S_n)=\Z$, $b_2(S_n)=n$, and elliptic axis of square $-n$.
For $s\ne0$, the surface $\mathcal K_s$ is obtained from the diagonal
Hopf surface $H_{\kappa(s)s^n,a(s)}$ by $n$ successive blow-ups at one
axis point and its indicated infinitely near strict-axis points.
In a fixed annulus marking of $\mathcal E$, its normal line has factor
\[
 g_s(v)=\kappa(s)(a(s)v+s)^n.
\]
This line extends across $s=0$ and is isomorphic to the holomorphic
family with factor $\kappa(s)a(s)^n v^n$.
For constant $a,\kappa$ it is a constant family of degree $-n$ lines.

The relation with the normal line of the unmodified Hopf axis is
\[
 N_{\mathcal E_s/\mathcal K_s}
   =L_{\kappa(s)s^n}\otimes\OO_{\mathcal E_s}(-n[-s]),
       \qquad s\ne0.
\]
Here $n[-s]$ is the weighted total exceptional correction, not the
intersection with the reduced final exceptional locus. It has no
Cartier extension on $\mathcal E$, even after a finite base change or
the addition of a vertical component. Thus this construction does
not contradict B.4.8: the marked divisor used to change the normal
line fails to extend.

For $n=8$, the degree $-8$ axis normal can be chosen to be the inverse
of $\OO_E(2)$ on any chosen smooth elliptic complete intersection
$E=(2,2)\subset\PP^3$. This supplies the normal-line matching with
$\operatorname{Bl}_E\PP^3$ described below, including one positive
triple section. It is an auxiliary construction and first-order
matching result. No attachment to the prescribed punctured
$S^6$ family, compact third component, or global smoothing is asserted.
\end{proposition}
\begin{proof}
We give the gluing domains as well as the local formulas. Let
$B_R\subset\C^2$ be a ball, and let
$\pi_n:\widehat B_R\to B_R$ be the following modification.
Blow up the origin, then successively the intersection of the newest
exceptional curve with the strict transform of the axis $x=0$.
In the last chart,
\[
 \pi_n(z,w)=(zw^n,w).
\]
Define the parameter-dependent embedding in that chart by
\[
 \sigma_s(z,w)=
       \bigl(\kappa(s)z,\ a(s)w+(1-a(s))s\bigr).
\]
Choose constants $A<\rho<1$ and $K>0$ with
$|a(s)|\le A$ and $|\kappa(s)|\le K$ on a closed smaller disc.
Shrink it further so that
$|a(s)w+(1-a(s))s|\le\rho R$ on $\overline B_R$.
Then
\[
 |\pi_n\sigma_s(z,w)|^2
 \le K^2R^2(\rho R)^{2n}+\rho^2R^2<R^2
\]
when $R$ is small enough. The last-chart coordinates of these images
are bounded, so this is compact containment of
$\sigma_s(\overline B_R)$ in $\widehat B_R$, uniformly in $s$.
Notice that no special smallness condition on the Picard parameter
$\kappa$ is required.

Remove the embedded ball and identify the outer and inner collars by
$\sigma_s\pi_n$. The outer collar avoids the exceptional locus, and
$\pi_n$ is a biholomorphism there; thus the collar identification is
a biholomorphism. Carry out the same construction over the parameter
disc. The embeddings and the identifications are holomorphic and
preserve $s$. Uniformly disjoint collars give a Hausdorff smooth
threefold, and the projection is a holomorphic submersion in every
chart. Uniform compact fundamental regions over compact parameter
sets prove properness. This is the usual Kato construction, with
the explicit moving embedding above; see
\citet[Section 2, pp.1351--1353]{IstratiEtAl2022Root}.
The general occurrence of deformations from Kato surfaces to
blown-up Hopf surfaces is classical; see
\citet[Introduction, pp.1750--1751]{Teleman2010AB}.
We use explicit formulas here to retain the elliptic marking and
its exceptional correction.

The contracted map $\pi_n\sigma_s$ fixes $(0,s)$.
In the coordinate $v=w-s$ it becomes
\[
 f_s(x,v)=
    \bigl(\kappa(s)x(a(s)v+s)^n,\ a(s)v\bigr).
\]
At $s=0$ this is the parabolic monomial germ with exponent matrix
$\left(\begin{smallmatrix}1&n\\0&1\end{smallmatrix}\right)$.
The chosen modification is exactly the $n$-fold modification in
\citet[Example 9.10, p.1386]{IstratiEtAl2022Root}.
Consequently the central surface is the stated parabolic Inoue
surface; its elliptic axis is disjoint from the rational cycle.
The construction gives $\pi_1=\Z$ and $b_2=n$, or equivalently one
can use the cited description of this surface. These are invariants
of the auxiliary surface, not of a proposed $S^6$ central fibre.

The strict axis is preserved in the collar gluing. Its coordinate
$v$ is multiplied by $a(s)$. To see its full quotient, the axis disc
in the blowdown coordinates is
\[
 D_s=\{v:|v+s|<R\}.
\]
For small $s$, it contains $0$ and satisfies $a(s)D_s\Subset D_s$.
Its dilates have union $\C$ and shrinking intersection $\{0\}$.
Thus $D_s\setminus a(s)D_s$, with the two boundary collars glued,
is a fundamental annulus for $\C^*/a(s)^\Z$.
This proves the exact elliptic marking throughout the parameter disc.

There are two normal frames, whose distinction is essential.
In the last blow-up chart the collar map $\sigma_s\pi_n$ has factor
\[
 g_z(v)=\kappa(s)(v+s)^n.
\]
On a fixed annulus avoiding $v=0$, shrink $s$ so that both
$v+s$ and $a(s)v+s$ are nonzero. The blowdown normal coordinate is
$x=z(v+s)^n$, and in that frame the factor is
\[
 g_x(v)=\kappa(s)(a(s)v+s)^n.
\]
The change between the frames is a nonzero gauge on these annuli.
The central factor is $\kappa(0)a(0)^n v^n$ and has degree $-n$,
with the factor-of-automorphy convention of
\citet[Proposition 5.6, pp.1369--1370]{IstratiEtAl2022Root}.

An explicit regular gauge verifies the whole family. On a fixed
annulus and its overlap with its $a(s)$-translate put
\[
 H_s(v)=\prod_{j\ge0}\left(1+\frac{a(s)^j s}{v}\right)^n.
\]
The product converges locally uniformly in $(v,s)$ and is nonzero
after shrinking the parameter disc. It satisfies
\[
 \frac{H_s(a(s)v)}{H_s(v)}
      =\left(1+\frac{s}{a(s)v}\right)^n.
\]
Replacing $x$ by $x/H_s(v)$ changes $g_x$ to
$\kappa(s)a(s)^n v^n$.
The collar maps are linear in the normal coordinate, so the same
calculation identifies an actual tubular line neighbourhood, not
only its first infinitesimal normal bundle.

For completeness, the generic Hopf blowdown also has an explicit
developing map. Fix $s\ne0$ and set
\[
 G_s(v)=\prod_{j\ge1}\left(1+\frac{a(s)^j v}{s}\right)^n,
       \qquad
 \Psi_s(x,v)=(xG_s(v),v).
\]
This product is entire in $v$, converges locally uniformly, and obeys
\[
 G_s(a(s)v)=\frac{G_s(v)}{(1+a(s)v/s)^n},\qquad
 \Psi_s f_s=\operatorname{diag}(\kappa(s)s^n,a(s))\Psi_s.
\]
For sufficiently small nonzero $s$, both diagonal multipliers have
absolute value less than one. Compose $\Psi_s$ with $\pi_n$ on a
Kato fundamental region. The displayed equivariance glues it to a
holomorphic map to $H_{\kappa(s)s^n,a(s)}$.
It never reaches the excluded origin: that would require
$x=v=0$, which is the fixed point $(0,s)$ removed with the inner ball.
All zero-factor curves of $G_s$ have
$v=-s/a(s)^j$, so they lie in the one axis orbit $[-s]$.
The original modification also contracts over $v=-s$.
At that representative,
\[
 G_s(-s)=\prod_{j\ge1}(1-a(s)^j)^n\ne0.
\]
Hence locally the map is precisely $\pi_n$ followed by a coordinate
unit. Away from that orbit it is a local biholomorphism of degree
one. Compactness makes it a proper modification, and the displayed
local model identifies it as the stated $n$ successive blow-ups.
This is a punctured-parameter construction: no meromorphic extension
of $G_s$ through $s=0$ is being claimed.

Each blow-up of a point on the strict elliptic curve subtracts that
point from its normal line. Following all $n$ stages gives
\[
 N_s=L_{\kappa(s)s^n}\otimes\OO_{\mathcal E_s}(-n[-s]).
\]
The multiplicity $n$ records total exceptional transforms.
The reduced final exceptional locus meets the strict axis only once.
One may check the formula directly: the standard theta factor
$\theta_a(av)=-v^{-1}\theta_a(v)$ shows that $\OO([p])$
has factor $-p/v$. At $p=-s$ the tensor product above has factor
\[
 \kappa(s)s^n(v/s)^n=\kappa(s)v^n,
\]
which is equivalent to $\kappa(s)a(s)^n v^n$ by a nonzero monomial
gauge. This checks the signs and degrees independently of the
successive-blow-up formula.

The divisor $n[-s]$ cannot extend. For constant $a$, choose two
distinct classes $[v_1],[v_2]\in\C^*/a^\Z$.
The sequences $s_j=-a^jv_i$ tend to zero, whereas the corresponding
divisors are constantly $n[v_i]$ and have different limits in
$\operatorname{Sym}^n(E_a)$.
For varying $a(s)$, choose the holomorphic logarithm
$\tau_a(s)=\log a(s)/(2\pi i)$ on the disc.
The Abel--Jacobi coordinate of $n[-s]$ is
$n\log(-s)/(2\pi i)$ modulo $\Z+\tau_a(s)\Z$.
On a positive parameter loop it winds $n$ times around a primitive
cycle in the smooth elliptic family. An extending section over the
disc would have zero winding. A finite base change of order $N$
multiplies this nonzero integer by $N$.
Adding a vertical elliptic component does not help: remove that
Cartier component, and the remaining proper horizontal divisor of
degree $n$ would be finite over the disc and define a continuous
section of the relative symmetric power, which has just been ruled
out. Nonproper exceptional divisors in deformations of Kato surfaces
are a classical phenomenon, studied by
\citet[Section 1.2]{DlousskyTeleman2012AB}; the calculation here
identifies the precise weighted divisor in this marked family.

This also explains the scope of ordinary elementary modifications.
If a degree-zero line $L_m$ from B.4.8 is changed by point blow-ups,
then its new normal is $L_m\otimes\OO(-D_{\rm exc})$.
If both that normal and the weighted Cartier correction extended on
a semistable good-reduction interface, their tensor product would
extend $L_m$, contradicting B.4.8. Drifting translations cannot
remove both requirements. The present Kato family escapes exactly
because $D_{\rm exc}=n[-s]$ has no such extension.

Finally let $E=(2,2)\subset\PP^3$ be smooth and set
$L=\OO_E(2)$, of degree eight. Its normal in $\PP^3$ is
$L\oplus L$. For $A=\operatorname{Bl}_E\PP^3$, the exceptional
divisor is $D=E\times\PP^1$ and, in the lines convention,
\[
 N_{D/A}=L\boxtimes\OO_{\PP^1}(-1).
\]
A bare Hopf-product cap has normal $M\boxtimes\OO$ with
$\deg M=0$. The required triple-divisor line would then be
$(L\otimes M)^{-1}\boxtimes\OO(1)$, which has no nonzero section
because its elliptic degree is $-8$. Thus no effective triple divisor
repairs that particular bare interface.

For the parabolic surface with $n=8$, changing $\kappa$ varies the
flat factor through all of $\operatorname{Pic}^0(E)$:
constant factors are parametrized by $\C^*/a^\Z$.
Choose its axis normal to equal $L^{-1}$.
In $B=S_8\times\PP^1$ the divisor $D=E\times\PP^1$ then has normal
$L^{-1}\boxtimes\OO$. With triple section $T=E\times\{p\}$,
\[
 N_{D/A}\otimes N_{D/B}\otimes\OO_D(T)
   =(L\boxtimes\OO(-1))\otimes
     (L^{-1}\boxtimes\OO)\otimes(\OO\boxtimes\OO(1))
   \simeq\OO_D.
\]
This proves precisely the first-order matching claimed.
\end{proof}

\begin{remark}[Auxiliary normal matching and the actual filling obstruction]
Choose a smooth quadric $Q\subset\PP^3$ containing $E$.
Its strict transform in $A$ is still $Q\simeq\PP^1\times\PP^1$:
blowing up its Cartier divisor $E$ does not change $Q$.
Its normal is
$\OO_Q(2,2)(-E)\simeq\OO_Q$, and it meets $D$ in a section.
The matching surface $S_8\times\{p\}$ in $B$ also has trivial
normal. A third threefold would therefore need divisors $Q$ and
$S_8$ meeting along $E$, with normals
$\OO_Q(-E)$ and $\OO_{S_8}(-E)$ respectively.
These are exactly the normal conditions for the central surface
$Q\cup_E S_8$ in a semistable surface smoothing.
The surface smoothing is constructed in B.4.11--B.4.12.
Its topological Euler number is $4+8-2\chi(E)=12$ and its
holomorphic Euler number is one. Its total space over an open
disc is not a compact third threefold, so that construction
does not supply a compact threefold assembly.

A full assembly must also preserve the required coherent cohomology
from B.4.5; a lone parabolic component attached to rational components
can lose its $H^1(\OO)$ through the elliptic overlap.
Additional conductor surfaces meeting $E$ would contribute further
triple divisors and must be included in the normal equation.
Nor is a local change of the cusp parameter automatically a permitted
global base change of the prescribed family.
These auxiliary matching statements do not produce a filling
of the actual family. B.4.33 proves that no such proper analytic
filling exists after any finite parameter cover.
\end{remark}
